\section{Implications for Safety Evaluation}

The benchmark is designed for practical diagnostic work. Model-policy teams can use it to find vague boundaries before those boundaries become grading failures. Red-teamers can use it to separate prompt-injection susceptibility from ordinary task failure. System-card authors can use it to expose fragile categories hidden by aggregate pass rates. External evaluators can use its provenance and claim register to ask exactly what a result licenses. These uses remain item-, procedure-, and configuration-specific unless their own validation and projection requirements are met (Section~\ref{sec:limitations}).

On the delegation-assurance account \parencite{reynolds2026delegationAssurance}, the benchmark supplies one measurement layer: evidence about whether model behaviour responds appropriately to instruction status, authority, scope, and failure attribution under controlled linguistic contrasts. It doesn't by itself establish institutional authorization, the adequacy of the resulting action record, or transport beyond the declared items, models, and wrappers; those require separate analyses. The companion evidentiary-assurance account \parencite{reynolds2026evidentiaryAssurance} treats the adequacy of that action record as its own layer. The three accounts divide the labour: this benchmark is the measurement layer, delegation assurance the normative-authorization layer, and evidentiary assurance the record-warrant layer.

The larger methodological point is that language-mediated distinctions are part of the control surface in these AI safety evaluations. A procedure that can't separate quotation from command, authority from content, policy analysis from unsafe enactment, or refusal from inability doesn't support the intended interpretation. A procedure that can make those separations on a fixed sample still needs separate evidence before its references, behaviours, measurements, or taxonomy are projected elsewhere.
